\begin{abstract}
We disprove the quadratic strengthening of the Sensitivity Conjecture by
constructing nonconstant total Boolean functions whose block sensitivity grows
faster than any constant multiple of sensitivity squared. In fact, for some
fixed \(\alpha>2\), our examples have unbounded block sensitivity and satisfy
\(\bs(f)\ge s(f)^\alpha\).
\end{abstract}
\section{Introduction}
\label{sec:introduction}

Sensitivity measures how many individual input bits can change a Boolean
function's value; block sensitivity allows disjoint sets of bits to be
flipped instead. We construct total Boolean functions for which the second
quantity grows faster than any constant multiple of the square of the first.

For a positive integer \(n\), write \([n]=\{1,\ldots,n\}\).
If \(x\in\{0,1\}^n\) and \(B\subseteq[n]\), let \(x^B\) be obtained
by flipping precisely the coordinates in \(B\); abbreviate
\(x^{\{a\}}\) to \(x^a\).
For a total Boolean function \(f:\{0,1\}^n\to\{0,1\}\), define
\[
 s(f,x)=\bigl|\{a\in[n]:f(x^a)\ne f(x)\}\bigr|,
 \qquad s(f)=\max_x s(f,x).
\]
Its block sensitivity \(\bs(f,x)\) is the largest number of pairwise
disjoint nonempty sets \(B_1,\ldots,B_m\subseteq[n]\) such that
\(f(x^{B_j})\ne f(x)\) for every \(j\), and
\(\bs(f)=\max_x\bs(f,x)\).

The Sensitivity Conjecture asked whether block sensitivity is bounded by a
polynomial in sensitivity, uniformly over the number of inputs and the
function. The question goes back to Nisan's work on parallel computation,
which introduced block sensitivity and already recorded Rubinstein's
quadratic separation \cite[Section~2]{Nisan1989}.
Nisan and Szegedy studied the problem through the degree of the unique
multilinear real polynomial representing a Boolean function, and explicitly
suggested the quadratic inequality \(\bs(f)\le s(f)^2\)
\cite[Section~4]{NisanSzegedy1994}.
Rubinstein's construction \cite{Rubinstein1995} showed that a polynomial
upper bound would need degree at least two. Virza improved the separation
to \(\bs(f)=\tfrac12s(f)^2+\tfrac12s(f)\) \cite{Virza2011}, and
Ambainis and Sun subsequently obtained examples satisfying
\[
 \bs(f)=\frac23s(f)^2-\frac13s(f)
\]
\cite[Theorem~1]{AmbainisSun2011}.
These improvements left open whether any universal quadratic upper bound
could hold.

Huang resolved the polynomial conjecture in 2019. His degree bound,
together with Tal's refinement of the degree--block-sensitivity comparison
\cite{Tal2013}, gives \(\bs(f)\le s(f)^4\) for every total Boolean function
\cite[Theorem~1.5]{Huang2019}. Wellens later improved its leading constant,
proving
\(\bs(f)\le\sqrt{2/3}\,s(f)^4+1\)
\cite[Corollary~4.2]{Wellens2022}.
These results settle the polynomial question. The theorem below rules out
its quadratic strengthening.

\begin{theorem}\label{thm:main}
For every real \(C>0\), there are a positive integer \(n\) and a
nonconstant total Boolean function \(f:\{0,1\}^n\to\{0,1\}\) such that
\[
 \bs(f)>C\,s(f)^2.
\]
More precisely, for every integer \(d\ge1\) the construction below gives
a nonconstant total Boolean function \(f\) with
\[
 \frac{\bs(f)}{s(f)^2}\ge \frac{2^d}{4(d+2)^2}.
\]
\end{theorem}
